	\documentclass[14pt]{extreport}
\usepackage{extsizes}
	\usepackage[frenchb]{babel}
	\usepackage[utf8]{inputenc}  
	\usepackage[T1]{fontenc}
	\usepackage{amssymb}
	\usepackage[mathscr]{euscript}
	\usepackage{stmaryrd}
	\usepackage{amsmath}
	\usepackage{tikz}
	\usepackage[all,cmtip]{xy}
	\usepackage{amsthm}
	\usepackage{varioref}
	\usepackage[ margin=.5in]{geometry}
	\geometry{a4paper}
	\usepackage{lmodern}
	\usepackage{hyperref}
	\usepackage{array}
	\usepackage{float}
	\usepackage{easytable}
	 \usepackage{fancyhdr}\usepackage{longtable}
	 \usetikzlibrary{shapes.misc}
\newlength{\taillecellule}
\setlength{\taillecellule}{2cm}
\newcolumntype{C}{@{}>{\centering\arraybackslash}p{\taillecellule}@{}}

\usepackage{pstricks,multido}
\usepackage{arrayjob}
\usepackage{calc,xlop}
\tikzset{cross/.style={cross out, draw=black, minimum size=2*(#1-\pgflinewidth), inner sep=0pt, outer sep=0pt},
%default radius will be 1pt. 
cross/.default={1pt}}

	\pagestyle{fancy}
	\theoremstyle{plain}
	\fancyfoot[C]{\empty} 
	\fancyhead[L]{Nom, prénom : }
	\fancyhead[R]{10 février 2026}
	
	
	\title{Interrogation chapitre 5}
	\date{}
	\begin{document}

\subsection*{Exercice 1 - 3 points}

Quelles consignes permettent de tracer un triangle ? Justifier votre réponse.
\begin{enumerate}
\item $AB = 5$ cm, $AC = 8$ cm, $BC = 2$ cm ;
\item $AB = 5$ cm, $AC = 3$ cm, $BC = 5$ cm ;
\item $AB = 4$ cm, $AC = 3$ cm, $BC = 5$ cm. 
\end{enumerate}


\subsection*{Exercice 2 - 4 points}

Calculez les sommes suivantes en détaillant vos étapes. 
\[ -5 - 1 - 2 =\]
\[ -1  + 3 + 2 + (-4) + (-5) + 6  = \] 
\[ (-1,5) + 3,5 + (+2,5) + (-4,5) + 5 \] 
\[ (-1,12) + 3,4 + (+2,12) + (-5,2) + (-0,2)\]

\subsection*{Exercice 3 - 4 points}
1) Indiquer les abscisses des points $B$, $C$, $D$ et $E$.
 
2) Placer sur la droite graduée le point $F$ d'abscisse $-1$. 

3) Quel point a une abscisse opposée de celle de $D$ ? 
 \begin{figure}[H]\center
\begin{tikzpicture}
\draw [->, >= latex] (-8, 0) -- (6, 0);
\draw (-7, -.2) -- (-7, .2);
\draw (-6, -.2) -- (-6, .2);
\draw (-5, -.2) -- (-5, .2);
\draw (-4, -.2) -- (-4, .2);
\draw(-3, -.2) -- (-3, .2);
\draw(-2, -.2) -- (-2, .2);
\draw (-1, -.2) -- (-1, .2);
\draw (5, -.2) -- (5, .2);
\draw (4, -.2) -- (4, .2);
\draw(3, -.2) -- (3, .2);
\draw(2, -.2) -- (2, .2);
\draw (1, -.2) -- (1, .2);
\draw (0, -.2) -- (0, .2);
\draw (0, .5) node {$O$};
\draw (0, -.5) node {$0$};
\draw (2, 0) node[cross=5pt, rotate = 0] {}; 
\draw (2, .5) node {$A$}; 
\draw (2, -.5) node {$1$}; 
\fill (-6, 0) node[cross=5pt, rotate = 0] {}; 
\draw (-6, .5) node {$B$}; 
\fill (5, 0) node[cross=5pt, rotate = 0] {};  
\draw (5, .5) node {$C$}; 
\fill (4, 0) node[cross=5pt, rotate = 0] {};  
\draw (4, .5) node {$D$}; 
\fill (-4, 0) node[cross=5pt, rotate = 0] {};  
\draw (-4, .5) node {$E$}; 
\end{tikzpicture}
\end{figure}
\subsection*{ Exercice 4 - 3 points}
Trier dans l'ordre croissant les nombres suivants : 
\[ -3,5 ;\phantom{miaou} 2 ;\phantom{miaou} 2,4 ;\phantom{miaou} -1,7 ;\phantom{miaou} -3,1 ;\phantom{miaou} 1,75 ;\phantom{miaou} 0 \]
\subsection*{Exercice 5 - 6 points}

Remplissez les pyramides additives suivantes (chaque nombre est la somme des deux nombres en-dessous de lui.)

\begin{figure}[H]
\center
\begin{tikzpicture}[scale=.9,every node/.style={draw,minimum width=1.8cm,minimum height=.9cm}]
\draw (0,0)  node {\phantom{9}};
\draw(-1,-1) node {\phantom {8}} ++(2,0) node {\phantom {1}};
\draw(-2,-2) node {\phantom{6,5}} ++(2,0) node {\phantom{1,5}} ++(2,0) node {\phantom{-0,5}};
\draw(-3,-3) node {\phantom{}{7,5}} ++(2,0) node {\phantom{}{-1}} ++(2,0) node {\phantom{}{2,5}} ++(2,0) node {\phantom{}{-3}};
\end{tikzpicture}
\
\begin{tikzpicture}[scale=.9,every node/.style={draw,minimum width=1.8cm,minimum height=.9cm}]
\draw (0,0)  node {\phantom{2,7}};
\draw(-1,-1) node {\phantom {1,2}} ++(2,0) node {\phantom {}{1,5}};
\draw(-2,-2) node {\phantom{1,8}} ++(2,0) node {\phantom{}{-0,6}} ++(2,0) node {\phantom{2,1}};
\draw(-3,-3) node {\phantom{}{4,1}} ++(2,0) node {\phantom{}{-2,3}} ++(2,0) node {\phantom{1,7}} ++(2,0) node {\phantom{0,4}};
\end{tikzpicture}
\end{figure}

 

 
\newpage
\subsection*{Exercice 1 - 3 points}
Quelles consignes permettent de tracer un triangle ? Justifier votre réponse.
\begin{enumerate}
\item $AB = 5$ cm, $AC = 8$ cm, $BC = 2$ cm ;
\item $AB = 5$ cm, $AC = 3$ cm, $BC = 5$ cm ;
\item $AB = 4$ cm, $AC = 3$ cm, $BC = 5$ cm. 
\end{enumerate}
\subsection*{Exercice 2 - 4 points}  
Calculez les sommes suivantes en détaillant vos étapes. 
\[ -4 - 3 - 6 = \]
\[ -1  + 4 + 2 + (-4) + (-9) + 6  = \] 
\[ (-4,5) + 1,5 + (+2,5) + (-4,5) + 5 \] 
\[ (-1,52) + 4,8 + (+2,52) + (-5,4) + (-0,4)\]

\subsection*{Exercice 3 - 4 points}
1) Indiquer les abscisses des points $B$, $C$, $D$ et $E$.
 
2) Placer sur la droite graduée le point $F$ d'abscisse $-2,5$. 

3) Quel point a une abscisse opposée de celle de $D$ ? 
 \begin{figure}[H]\center
\begin{tikzpicture}
\draw [->, >= latex] (-8, 0) -- (6, 0);
\draw (-7, -.2) -- (-7, .2);
\draw (-6, -.2) -- (-6, .2);
\draw (-5, -.2) -- (-5, .2);
\draw (-4, -.2) -- (-4, .2);
\draw(-3, -.2) -- (-3, .2);
\draw(-2, -.2) -- (-2, .2);
\draw (-1, -.2) -- (-1, .2);
\draw (5, -.2) -- (5, .2);
\draw (4, -.2) -- (4, .2);
\draw(3, -.2) -- (3, .2);
\draw(2, -.2) -- (2, .2);
\draw (1, -.2) -- (1, .2);
\draw (0, -.2) -- (0, .2);
\draw (0, .5) node {$O$};
\draw (0, -.5) node {$0$};
\draw (2, 0) node[cross=5pt, rotate = 0] {}; 
\draw (2, .5) node {$A$}; 
\draw (2, -.5) node {$1$}; 
\fill (-6, 0) node[cross=5pt, rotate = 0] {}; 
\draw (-6, .5) node {$B$}; 
\fill (5, 0) node[cross=5pt, rotate = 0] {};  
\draw (5, .5) node {$C$}; 
\fill (-2, 0) node[cross=5pt, rotate = 0] {};  
\draw (-2, .5) node {$D$}; 
\fill (+3, 0) node[cross=5pt, rotate = 0] {};  
\draw (+3, .5) node {$E$}; 
\end{tikzpicture}
\end{figure}
\subsection*{ Exercice 4 - 3 points}
Trier dans l'ordre croissant les nombres suivants : 
\[ -3,5 ;\phantom{miaou} -2 ;\phantom{miaou} 2,4 ;\phantom{miaou} 1,7 ;\phantom{miaou} -3,1 ;\phantom{miaou} -1,75 ;\phantom{miaou} 0 \]

\subsection*{Exercice 5 - 6 points}

Remplissez les pyramides additives suivantes (chaque nombre est la somme des deux nombres en-dessous de lui.)

\begin{figure}[H]
\center
\begin{tikzpicture}[scale=.9,every node/.style={draw,minimum width=1.8cm,minimum height=.9cm}]
\draw (0,0)  node {\phantom{9}};
\draw(-1,-1) node {\phantom {8}} ++(2,0) node {\phantom {1}};
\draw(-2,-2) node {\phantom{6,5}} ++(2,0) node {\phantom{1,5}} ++(2,0) node {\phantom{-0,5}};
\draw(-3,-3) node {\phantom{}{7,5}} ++(2,0) node {\phantom{}{-1}} ++(2,0) node {\phantom{}{2,5}} ++(2,0) node {\phantom{}{-3}};
\end{tikzpicture}
\ 
\begin{tikzpicture}[scale=.9,every node/.style={draw,minimum width=1.8cm,minimum height=.9cm}]
\draw (0,0)  node {\phantom{0}};
\draw(-1,-1) node {\phantom{} {-5,9}} ++(2,0) node {\phantom {5,9}};
\draw(-2,-2) node {\phantom{-5,7}} ++(2,0) node {\phantom{}{-0,2}} ++(2,0) node {\phantom{6,1}};
\draw(-3,-3) node {\phantom{}{2,4}} ++(2,0) node {\phantom{-8,1}} ++(2,0) node {\phantom{7,9}} ++(2,0) node {\phantom{}{-1,8}};
\end{tikzpicture}
\end{figure} 
 
\end{document}